\section{Tensora Design}
\label{sec:design}

Tensora is the measurement framework for this paper. Its design goal is to make workload shape, backend choice, and integration overhead explicit enough to measure across I/O backends.

\subsection{Design principles}

A format parser should produce the same read plan regardless of backend, so timing differences reflect submission mechanism rather than work performed. Backend labels are explicit at the harness level. Measurement (cache drops, timing, capture) stays outside the loader.

\subsection{Architecture}

Tensora separates format parsing, backend execution, and evaluation (Figure~\ref{fig:tensora-arch}). The \textbf{format layer} parses checkpoint metadata and produces a request-local read plan: files, byte ranges, tensor names, and expected sizes. The \textbf{backend layer} executes reads through an explicit backend (sync, async, or the adaptive default) selected by the caller. The \textbf{evaluation layer} contains the Rust profiler and vLLM experiments; Python bindings bridge the integration path.

\begin{figure}[t]
  \centering
  \resizebox{\linewidth}{!}{%
  \begin{tikzpicture}[
    font=\small,
    box/.style={draw, rounded corners=2pt, align=center, inner sep=6pt, minimum width=2.6cm},
    arr/.style={->, >=latex, thick},
  ]
    \node[box] (fmt) {Format layer\\metadata $\rightarrow$ read plan};
    \node[box, right=1.1cm of fmt] (be) {Backend\\sync / async / default};
    \node[box, right=1.1cm of be] (out) {Tensor buffers\\Rust / Python FFI};
    \node[box, below=1.0cm of be] (ev) {Evaluation harness\\timing + cache state};
    \draw[arr] (fmt) -- (be);
    \draw[arr] (be) -- (out);
    \draw[arr, dashed] (ev) -- (be);
  \end{tikzpicture}%
  }
  \caption{Tensora architecture. The hot path is format parsing $\rightarrow$ backend execution $\rightarrow$ tensor buffers. The harness controls timing and cache state outside that path.}
  \label{fig:tensora-arch}
\end{figure}

All backends receive the same read plan from the format layer and execute reads through the selected I/O submission method.

\subsection{Common backend surface}

All eager backends implement the same conceptual operations:

\begin{itemize}
\item \texttt{load(path)} for one full file,
\item \texttt{load\_batch(paths)} for multiple whole-file loads,
\item \texttt{load\_range(path, offset, len)} for one byte range, and
\item \texttt{load\_range\_batch(requests)} for grouped range reads.
\end{itemize}

This surface is deliberately small. SafeTensors mostly exercises whole-file and batch loading; ServerlessLLM exercises range batches. Backend implementations may chunk, group, coalesce, or tune queue depth internally, but they must satisfy the same logical tensor plan.

\subsection{Backend implementations}

The \textbf{sync} backend uses blocking POSIX reads with practical host-side parallelism for large files. It is not a strawman serial baseline; its performance on small SafeTensors loads reflects real low-overhead bulk I/O where orchestration cost is minimal.

The \textbf{async} backend is built on Tokio. It groups work per file, batches dynamically, and avoids unnecessary task fanout. Its concurrent scheduling model amortizes per-request overhead when many operations are in flight, which becomes advantageous for larger SafeTensors workloads and for ServerlessLLM's range-dominated access pattern.

The \textbf{default} backend is an adaptive selector that chooses between sync and async (or \texttt{io\_uring} when available) based on checkpoint layout, model scale, and platform capability. It is designed to match or approach the best explicit backend at each point in the format$\times$scale space without requiring manual backend selection.

On platforms where \texttt{io\_uring} is available (bare-metal Linux with kernel $\geq$ 5.6), Tensora also exposes an \textbf{\texttt{io\_uring}} backend with persistent rings, batched submissions, and workload-dependent chunk planning. This backend is unavailable on Modal's gVisor containers (\texttt{ENOSYS}) and is therefore not part of the primary evaluation.

The mmap-backed path remains useful for lazy selective access, especially through Python handles, but it is not the main eager cold-start mechanism evaluated in the Rust matrices.

\subsection{Formats and integrations}

Tensora supports two layouts because they expose different regimes. \texttt{SafeTensors} presents whole-file or multi-shard contiguous loading. \texttt{ServerlessLLM} presents an indexed partition layout with range-oriented reads. Supporting both under one framework is what makes the layout-dependent regime claim measurable.

Python bindings are implemented with PyO3 and expose a small API: sync and async eager load, and lazy \texttt{open\_*} handles. The async Python path uses a shared Tokio runtime rather than constructing one per call. The vLLM integration uses a custom loader path and follows vLLM's synchronous loader boundary, so the serving benchmark reflects the interface available to real integrations rather than an invented async serving path. The Python installation is the bridge layer through which vLLM accesses Tensora.
